
\begin{frame}[fragile]{Algèbre relationnelle\,: les opérateurs}

SQL propose un autre type d'interrogation, \alert{fonctionnelle},
basée sur \alert{l'algèbre relationnelle}. 

\vfill

L'algèbre est un ensemble de 6 opérateurs, qui présentent deux propriétés essentielles

\begin{redbullet}
  \item \alert{Clôture}: un opérateur s'applique à des relations et produit une relation
  \item \alert{Composition}: un opérateur peut prendre en entrée le résultat d'un autre pour
      définir des \alert{requêtes} algébriques complexes.
\end{redbullet}

\vfill

Cette session présente les 6 opérateurs nécessaires et suffisants.

\vfill
\begin{blocgauche}
\alert{Ces diapositives correspondent au support en ligne disponible sur
le site \url{http://sql.bdpedia.fr/}}
\end{blocgauche}

\end{frame}


\begin{frame}[fragile]{La projection }

$\pi_{A_1, A_2, \ldots,A_k}(R)$ construit une relation contenant 
tous les nuplets de $R$,  dans lesquels seuls les attributs $A_1, A_2, \ldots A_k$ 
sont conservés. 

\vfill

Exemple\,: nom et lieu des logements: $\pi_{nom, lieu}(Logement) $

\vfill

\begin{minted}[baselinestretch=.9,
               fontsize=\small,
               framesep=2mm]{sql}
select nom, lieu
from Logement
\end{minted}
\vfill
\begin{tabular}{c|c}
\textbf{nom} & \textbf{lieu}\\ \hline 
Causses & Cévennes\\ \hline 
Génépi & Alpes\\ \hline 
U Pinzutu & Corse\\ \hline 
Tabriz & Bretagne\\ \hline 
\end{tabular}

\vfill

\begin{blocgauche}
\sqlkw{select *} permet de conserver tous les attributs.
\end{blocgauche}
\end{frame}

%%%%%%%%%%%%%%%%%%%%
\begin{frame}[fragile]{La sélection, $\sigma$}

La sélection $\sigma_F(R)$ s'applique à une relation, $R$, et
en extrait les nuplets qui satisfont $F$
$$\sigma_{lieu='Corse'}(Logement)$$
\vfill
En SQL: 

\begin{minted}[baselinestretch=.9,
               fontsize=\small,
               framesep=2mm]{sql}
select * 
from Logement
where lieu = 'Corse'
\end{minted}

\vfill

\begin{blocgauche}
Les comparaisons s'écrivent
$A \Theta B$, où $\Theta$ appartient
à $\{=, <, >, \leq, \geq\}$.
\end{blocgauche}
\end{frame}


%%%%%%%%%%%%%%%%%%%%
\begin{frame}[fragile]{Le produit cartésien, $\times$}

$R \times S$ produit 
une relation où chaque nuplet de $R$ est associé
à chaque nuplet de $S$.

\vfill

\begin{minipage}[t]{6cm}
\begin{center}
\begin{tabular}{c|c}
A & B \\ \hline
a & b  \\
x & y  \\\hline 
\end{tabular}
\end{center}
\end{minipage} \
\begin{minipage}[t]{6cm}
\begin{center}
\begin{tabular}{c|c}
C & D \\ \hline
c & d  \\ 
u & v  \\
x & y  \\\hline 
\end{tabular}
\end{center}
\end{minipage} 

\vfill


Et voici le résultat de $R \times S$\,:

\begin{center}
\begin{tabular}{c|c|c|c}
A & B & C & D \\ \hline
a & b & c & d  \\ 
a & b & u & v \\
a & b & x & y \\
x & y & c & d  \\ 
x & y & u & v \\
x & y & x & y \\\hline
\end{tabular}
\end{center}

\end{frame}


%%%%%%%%%%%%%%%%%%%%
\begin{frame}[fragile]{En SQL, \sqlkw{cross join}}

En SQL, le produit cartésien s'exprime avec \texttt{cross join}. 
On place l'expression algébrique \alert{dans} la clause \sqlkw{from}.

\vfill


\begin{minted}[baselinestretch=.9,
               fontsize=\small,
               framesep=2mm]{sql}
select * 
from R cross join S
\end{minted}

\vfill

\alert{Raisonnement}: \texttt{R \sqlkw{cross join} S} définit une  nouvelle relation, les 
autres opérateurs (sélection, projection) s'appliquent à cette relation.

\vfill

\begin{blocgauche}
Le \sqlkw{from} définit donc une relation \alert{calculée}
\end{blocgauche}
\end{frame}

%%%%%%%%%%%%%%%%%%%%
\begin{frame}[fragile]{Problème des noms d'attribut ambigus}

Il peut arriver que les deux relations
aient des attributs qui ont le même nom. 

\vfill

On doit alors  donner un nom distinct à chaque attribut.

\vfill

\begin{minipage}[t]{6cm}
\begin{center}
Relation $R$:
\begin{tabular}{c|c}
A & B \\ \hline
a & b  \\
x & y  \\\hline 
\end{tabular}
\end{center}
\end{minipage} \
\begin{minipage}[t]{6cm}
\begin{center}
Relation $T$:
\begin{tabular}{c|c}
A & B \\ \hline
m & n  \\
o & p  \\\hline 
\end{tabular}
\end{center}
\end{minipage} 


\vfill
\begin{blocgauche}
$R \times T$
a pour schéma $(A, B, A, B)$ et présente
donc des ambiguités,
\end{blocgauche}
\end{frame}


%%%%%%%%%%%%%%%%%%%%
\begin{frame}[fragile]{Renommage $\rho$}

L'expression $\rho_{A \to C,B \to D}(T)$ renomme 
$A$ en $C$ et $B$ en $D$ dans la relation
$T$.

\vfill

$\rho_{A\to C,B \to D}(T)$:
\begin{tabular}{c|c}
C & D \\ \hline
m & n  \\
o & p  \\\hline 
\end{tabular}

\vfill

Le résultat de $R \times \rho_{A\to C,B \to D}(T)$ 

\vfill

\begin{tabular}{c|c|c|c}
A & B & C & D \\ \hline
a & b & m & n  \\ 
a & b & o & p \\
x & y & m & n  \\ 
x & y & o & p \\\hline
\end{tabular}

\end{frame}



%%%%%%%%%%%%%%%%%%%%
\begin{frame}[fragile]{Renommage en SQL: \sqlkw{as}}


La requête
\
\vfill

\begin{minted}[baselinestretch=.9,
               fontsize=\small,
               framesep=2mm]{sql}
select Voyageur.idVoyageur, Séjour.idVoyageur
from Voyageur cross join Séjour
\end{minted}

engendre une erreur \textit{duplicate field name:}

\vfill

Bonne version:

\begin{minted}[baselinestretch=.9,
               fontsize=\small,
               framesep=2mm]{sql}
select Voyageur.idVoyageur as idV1, Séjour.idVoyageur as idV2
from Voyageur cross join Séjour
\end{minted}

\begin{blocgauche}
Le \sqlkw{as} permet aussi de renommer des relations. Cf. support de cours.
\end{blocgauche}

\end{frame}


%%%%%%%%%%%%%%%%%%%%
\begin{frame}[fragile]{L'union, $\cup$}

$R \cup S$ produit 
une relation contenant l'union de $R$ $S$ (qui doivent avoir le même schéma).

\vfill

\begin{minipage}[t]{6cm}
\begin{center}
\begin{tabular}{c|c}
A & B \\ \hline
a & b  \\
x & y  \\\hline 
\end{tabular}
\end{center}
\end{minipage} \
\begin{minipage}[t]{6cm}
\begin{center}
\begin{tabular}{c|c}
A & B\\ \hline
c & d  \\ 
u & v  \\
x & y  \\\hline 
\end{tabular}
\end{center}
\end{minipage} 

\vfill


Et voici le résultat de $R \cup S$\,:

\begin{center}
\begin{tabular}{c|c}
A & B\\ \hline
a & b  \\
c & d  \\ 
u & v  \\
x & y  \\\hline 
\end{tabular}
\end{center}

\end{frame}


%%%%%%%%%%%%%%%%%%%%
\begin{frame}[fragile]{Union en SQL: \sqlkw{union}}


La requête donnant l'union des noms de lieu et des noms de logement
\
\vfill

\begin{minted}[baselinestretch=.9,
               fontsize=\small,
               framesep=2mm]{sql}
     select lieu from Logement
       union
     select région as lieu from Voyageur
\end{minted}

\vfill

Nb\,: le schéma du résultat est \texttt{(lieu)}

\vfill

\begin{blocgauche}
Rarement utilisée, mais indispensable (pas d'autre expression possible)
en cas de besoin
\end{blocgauche}

\end{frame}


%%%%%%%%%%%%%%%%%%%%
\begin{frame}[fragile]{La différence, $-$}

$R \- S$ produit 
une relation contenant les nuplets de $R$ qui ne sont pas
dans $S$ (elles doivent avoir le même schéma).

\vfill

\begin{minipage}[t]{4cm}
\begin{center}
Relation $R$
\begin{tabular}{c|c}
A & B \\ \hline
a & b  \\
x & y  \\\hline 
\end{tabular}

\end{center}
\end{minipage} \
\begin{minipage}[t]{4cm}
\begin{center}
Relation $S$
\begin{tabular}{c|c}
A & B\\ \hline
c & d  \\ 
u & v  \\
x & y  \\\hline 
\end{tabular}
\end{center}
\end{minipage} 
\
\begin{minipage}[t]{4cm}
\begin{center}
Relation $R - S$
\begin{tabular}{c|c}
A & B\\ \hline
a & b  \\\hline 
\end{tabular}

\end{center}
\end{minipage} 
\vfill

\vfill 
En SQL


\begin{minted}[baselinestretch=.9,
               fontsize=\small,
               framesep=2mm]{sql}
     select lieu from Logement
       except
     select région as lieu from Voyageur
\end{minted}

\vfill
\begin{blocgauche}
Très peu pratique à cause de la contrainte sur les schémas. La
version déclarative, \sqlkw{not exists}, est bien plus facile.
\end{blocgauche}

\end{frame}

%%%%%%%%%%%
\begin{frame}{À retenir}


\vfill

\begin{redbullet}
   \item Syntaxe et signification des 6 opérateurs 
  \item  Notion de \alert{clôture}: toute opération s'applique à des relations
     et produit une relation
  \item Notion de \alert{composition}: on crée des requêtes complexes en combinant
  des opérateurs, pas en créant des opérateurs super-complexes.
  \item En SQL: l'algèbre opère sur des relations complètes, alors que dans la perspective
     déclarative, SQL exprime des conditions sur des nuplets.
\end{redbullet}

\vfill
\begin{blocgauche}
Nouveauté principale pour SQL: la possibilité d'introduire des expressions dans le \sqlkw{from}
\end{blocgauche}
\end{frame}

